\documentclass[conference]{IEEEtran}
\IEEEoverridecommandlockouts

% ── Packages ──────────────────────────────────────────────────────────────────
\usepackage{cite}
\usepackage{amsmath,amssymb,amsfonts}
\usepackage{algorithmic}
\usepackage{graphicx}
\usepackage{textcomp}
\usepackage{xcolor}
\usepackage{booktabs}
\usepackage{multirow}
\usepackage{array}
\usepackage{listings}
\usepackage{float}
\usepackage{hyperref}
\usepackage{url}
\usepackage{colortbl}
\usepackage{microtype}
\usepackage{balance}

% ── Colours ───────────────────────────────────────────────────────────────────
\definecolor{codebg}{RGB}{248,248,248}
\definecolor{coderule}{RGB}{200,200,200}
\definecolor{highlight}{RGB}{214,245,214}
\definecolor{accentblue}{RGB}{0,82,155}

% ── Listings (code blocks) ────────────────────────────────────────────────────
\lstset{
    backgroundcolor=\color{codebg},
    basicstyle=\ttfamily\scriptsize,
    breaklines=true,
    breakatwhitespace=false,
    columns=fullflexible,
    keepspaces=true,
    frame=single,
    rulecolor=\color{coderule},
    framesep=3pt,
    numbers=left,
    numberstyle=\tiny\color{gray},
    numbersep=4pt,
    xleftmargin=8pt,
    xrightmargin=0pt,
    captionpos=b,
    belowskip=4pt,
    aboveskip=4pt,
}

% ── Hyperref ──────────────────────────────────────────────────────────────────
\hypersetup{
    colorlinks=true,
    linkcolor=accentblue,
    filecolor=accentblue,
    urlcolor=accentblue,
    citecolor=accentblue
}

\begin{document}

% ── Title ─────────────────────────────────────────────────────────────────────
\title{%
    \textbf{AI Principal Tracker Ethos:}\\
    A Practical, Longitudinal System for\\
    Comparative AI Policy Analysis and Ethics Signals (2019--2025)%
}

\author{
    \IEEEauthorblockN{Souda Fathima Kurnool}
    \IEEEauthorblockA{\textit{MSc in Artificial Intelligence}\\
    National College of Ireland, Dublin\\
    \href{mailto:24251445@student.ncirl.ie}{\texttt{24251445@student.ncirl.ie}}}
    \and
    \IEEEauthorblockN{Ali Akarsu}
    \IEEEauthorblockA{\textit{MSc in Artificial Intelligence}\\
    National College of Ireland, Dublin\\
    \href{mailto:24304051@student.ncirl.ie}{\texttt{24304051@student.ncirl.ie}}}
    \and
    \IEEEauthorblockN{Jenil Jayeshbhai Parmar}
    \IEEEauthorblockA{\textit{MSc in Artificial Intelligence}\\
    National College of Ireland, Dublin\\
    \href{mailto:24281883@student.ncirl.ie}{\texttt{24281883@student.ncirl.ie}}}
    \and
    \IEEEauthorblockN{Smit Rajendrasinh Baghel}
    \IEEEauthorblockA{\textit{MSc in Artificial Intelligence}\\
    National College of Ireland, Dublin\\
    \href{mailto:25117149@student.ncirl.ie}{\texttt{25117149@student.ncirl.ie}}}
    \and
    \IEEEauthorblockN{Dhruv Sharma}
    \IEEEauthorblockA{\textit{MSc in Artificial Intelligence}\\
    National College of Ireland, Dublin\\
    \href{mailto:24234915@student.ncirl.ie}{\texttt{24234915@student.ncirl.ie}}}
}

\maketitle

% ─────────────────────────────────────────────────────────────────────────────
\begin{abstract}
This paper presents the current implementation of \textbf{AI Principal Tracker
Ethos}, a full-stack web system for comparative, longitudinal AI policy
analysis. The project objective is to make corporate AI policy evolution
interpretable to non-experts by combining structured policy timelines with a
large-scale ethics keyword corpus. Two complementary data sources are
integrated: (1) structured policy timelines for Google, Microsoft, IBM, Amazon,
Meta, Apple, and Tesla (\textbf{537} records), and (2) a multi-year Gemini
ethics keyword corpus spanning 2019--2025 (\textbf{6,\,967} records) used to
derive category-level ethics emphasis (transparency, privacy, accountability,
safety, governance, and related signals). The system combines a FastAPI backend,
React + TypeScript frontend, and a lightweight LLM integration (Azure OpenAI
optional; mock default) to provide policy timelines, ethics summaries, and user
ratings. Results show that a keyword-based ethics timeline can provide a
transparent, reproducible proxy for longitudinal emphasis while retaining a
practical, fast, and deployable architecture.

\smallskip\noindent\textbf{Key Contributions:}
\begin{enumerate}
  \item A production-ready AI policy dashboard with professional timeline
    visualisations and category-level signals derived from real datasets.
  \item A Gemini-specific ethics timeline built from \textbf{6,\,967} keyword
    records across 2019--2025, enabling year-by-year ethics emphasis.
  \item A modular, low-cost architecture that supports mock analysis by
    default and optional Azure OpenAI summarisation without heavy local
    inference.
\end{enumerate}
\end{abstract}

% ─────────────────────────────────────────────────────────────────────────────
\section{Introduction}

AI policy documents are increasingly visible, yet users struggle to interpret
what they imply in practice. Prior work highlights the gap between stated
principles and operational reality---often labeled \emph{ethics washing} or
symbolic compliance~\cite{munn2023ethics,stanford2023hai}. Our system addresses
this by providing a concrete, user-facing tool to examine AI policy timelines
and ethics emphasis over time. Instead of claiming new global statistics, we
focus on transparent data processing, reproducible outputs, and clear
visualisation of policy signals for different companies.

\subsection*{Project Objectives and Novelty}
The project is novel in its practical framing: rather than building an abstract
ethics model, it integrates two datasets with different structure and granularity
to produce interpretable, comparative timelines. The novelty lies in (i)
repurposing a large ethics keyword corpus into a longitudinal ethics signal for
Gemini, (ii) combining it with explicit policy timelines for other companies,
and (iii) delivering this through a deployable dashboard that emphasises
transparency and explainability. This approach avoids heavy infrastructure while
still enabling longitudinal insights, aligning with the assessment emphasis on
originality and impact.

\subsection*{Problem Significance}
Major providers publish AI principles, but end users rarely see how these
principles evolve. Policy changes are subtle and dispersed across years, while
ethics clauses can remain ambiguous. This motivates a dashboard that:
(1) shows multi-year policy changes, (2) highlights ethics categories, and
(3) supports rapid comparison between companies and datasets. The outcome is a
tool that can support researchers, policy analysts, and end users who require
clear evidence of ethics trends and potential gaps.

\subsection*{Research Questions}
\begin{description}
  \item[RQ1] How can AI policy timelines be made interpretable and comparable
    across multiple companies and time periods?
  \item[RQ2] Can keyword-based ethics signals provide a useful longitudinal
    proxy for policy emphasis when direct policy text is unavailable?
  \item[RQ3] How can a lightweight LLM integration support policy summarisation
    without heavy infrastructure or high compute costs?
\end{description}

% ─────────────────────────────────────────────────────────────────────────────
\section{Related Work}

Existing work provides valuable ethics frameworks but limited empirical tooling.
Jobin et al. catalog a global inventory of ethics guidelines without a
longitudinal, user-facing system~\cite{jobin2019global}. Munn provides a
critical perspective on ethics washing but does not operationalise a pipeline
or artefact for public use~\cite{munn2023ethics}. The AI life-cycle guidelines
propose normative frameworks but are not grounded in concrete corporate policy
timelines~\cite{springer2022development}.

Our contribution is positioned differently: it emphasises a reproducible
pipeline, deployable interface, and practical dataset integration. By combining
structured policy timelines with an ethics keyword corpus and exposing the
results via an accessible UI, we move from critique to practical observation
and comparison. This aligns with calls for applied transparency and improved
public access to policy evolution evidence.

% ─────────────────────────────────────────────────────────────────────────────
\section{Methodology}

\subsection*{Data Sources}
Two datasets are integrated.

\begin{description}
  \item[Dataset A -- Corporate policy timelines] CSV-based policy points for
  major companies (Google, Microsoft, IBM, Amazon, Meta, Apple, Tesla) with
  \textbf{537} total records (256 Google; 281 Microsoft/IBM/Amazon).
  \item[Dataset B -- Gemini ethics keyword corpus] Five CSVs (2019, 2020,
  2022, 2024, 2025) totalling \textbf{6,\,967} keyword records.
  Each record maps a word to ethics categories such as fairness,
  transparency, privacy, accountability, safety, governance, and oversight.
\end{description}

\subsection*{Data Preprocessing and Storage}
Corporate policy CSVs are normalised into a common schema
\{company, year, policy\_point, category, severity, impact, status\}. The Gemini
corpus is parsed by year and mapped to ethics categories. We apply case
normalisation, whitespace trimming, and unique-word counting to avoid inflated
counts from repeated terms. Parsed records are cached in-memory at the backend
service layer to reduce repeated CSV reads.

\subsection*{Dataset Integration Strategy}
The project integrates heterogeneous data by aligning both datasets to a
common timeline interface. Corporate policy entries provide explicit
year-stamped policy points, while Gemini keywords are aggregated into ethics
signals per year. Both are mapped to shared fields (year, category, severity,
impact) so that visualisation and comparisons are consistent regardless of
source. This integration strategy is a core novelty of the system.

\subsection*{Pipeline Overview}
Figure~\ref{fig:pipeline} summarises the end-to-end workflow from CSV ingestion
to the UI timeline visualisation.

\begin{figure}[!htb]
  \centering
  \includegraphics[width=\columnwidth]{pipeline_diagram.png}
  \caption{System pipeline: CSV ingestion $\to$ ethics/category aggregation
    $\to$ timeline API $\to$ React visualisation and policy cards.}
  \label{fig:pipeline}
\end{figure}

\subsection*{Gemini Ethics Timeline Construction}
For each year, unique keywords are counted per ethics category to produce an
annual emphasis signal. The timeline visualises these counts as policy points
(e.g., \emph{Transparency Focus}) with an impact indicator that reports the
keyword count for that year. This provides a transparent, reproducible proxy
for ethics emphasis.

\subsection*{Lightweight ML Baseline (Semantic Drift)}
To strengthen model-based evaluation without heavy compute, we implement a
TF--IDF baseline on policy points per year and compute cosine similarity
between consecutive years. This provides a lightweight semantic drift signal
that complements the timeline view. The baseline is explainable, fast, and
aligned with the project's emphasis on interpretability.

\subsection*{System Architecture}
The application is implemented as a lightweight full-stack system:

\begin{itemize}
  \item \textbf{Backend:} FastAPI with CSV ingestion, policy timeline API, and
  optional LLM summarisation via Azure OpenAI.
  \item \textbf{Frontend:} React + TypeScript + Vite with a timeline UI,
  policy cards, quick review view, and user ratings interface.
  \item \textbf{Data:} Local CSV datasets and JSON-based ratings storage for
  fast setup and reproducibility.
\end{itemize}

\subsection*{User Review Framework}
The system includes a human-in-the-loop review layer in which users rate
company policies across fairness, transparency, privacy, and accountability.
Each rating is stored and aggregated per company. The \emph{Quick Review}
Dashboard visualises these aggregates, presenting both overall performance
and category-level breakdowns. This ensures that policy evaluation is grounded
in user judgment rather than only automated signals.

\subsection*{API Endpoints}
Key endpoints include:
\begin{itemize}
  \item \texttt{GET /ethics/companies} --- list available companies (including
  Gemini).
  \item \texttt{GET /ethics/timeline/\{company\}} --- policy timeline or Gemini
  ethics signals.
  \item \texttt{POST /chat} --- optional LLM-backed explanations (mock by
  default).
\end{itemize}

\subsection*{Challenges and Design Decisions}
Three implementation challenges informed the final design:
\begin{enumerate}
  \item \textbf{Heterogeneous data} --- corporate policy points are structured,
  while Gemini data is keyword-based. We designed a common timeline interface
  so both datasets can be rendered consistently.
  \item \textbf{Compute constraints} --- to avoid heavy local inference, the
  system defaults to a mock LLM provider with optional Azure OpenAI for
  analysis. This keeps the system deployable on modest hardware.
  \item \textbf{Interpretability} --- raw keyword counts are not intuitive, so
  we map them to category-level ``focus'' items with severity bands (low,
  medium, high) based on year-specific counts.
\end{enumerate}

\subsection*{Evaluation Strategy}
Evaluation focuses on data coverage, pipeline correctness, and UI clarity.
We report dataset sizes, year coverage, and qualitative inspection of timeline
rendering to ensure each year contains meaningful policy signals. The aim is
to provide a reproducible, interpretable artefact rather than opaque metrics.

% ─────────────────────────────────────────────────────────────────────────────
\section{Results and Evaluation}

\subsection*{Timeline Coverage}
The Gemini timeline spans 2019--2025 and exposes shifts across 20+ ethics
categories. The corporate timeline dataset provides explicit policy points
with category and severity fields, enabling cross-company comparisons.

\begin{table}[!htb]
  \renewcommand{\arraystretch}{1.1}
  \caption{Dataset coverage used in the current implementation.}
  \label{tab:datasets}
  \begin{tabular}{lrr}
    	oprule
    	extbf{Dataset} & \textbf{Records} & \textbf{Years} \\
    \midrule
    Gemini ethics keyword corpus & 6{,}967 & 2019--2025 \\
    Corporate policy timelines & 537 & 2018--2022 \\
    \bottomrule
  \end{tabular}
\end{table}

\subsection*{System Usability}
The updated Analyse section presents clear entry points for timeline analysis
and supports fast loading via cached CSV parsing. Policy cards display
category, severity, and impact in a professional, readable layout.

\subsection*{User Review Outcomes}
The Quick Review Dashboard consolidates user ratings into an overall score and
category-level metrics, giving stakeholders a concise view of company
performance alongside the timeline evidence. This user-driven layer enhances
the practical impact of the system by aligning policy evaluation with human
judgment.

\subsection*{Interpretation and Impact}
The Gemini timeline makes year-by-year ethics emphasis visible through a
consistent visual language. When combined with explicit corporate policy
points, the system enables a cross-dataset comparison of ethics commitments
and their evolution. This directly supports stakeholders who need interpretable
policy evidence rather than opaque summaries.

\subsection*{Semantic Drift Baseline}
The TF--IDF cosine similarity baseline provides a compact measure of policy
shift between consecutive years. While not a full predictive model, it offers
an interpretable signal for whether policy language is stabilising or changing,
and it can be extended to compare drift across companies in future work.

% ─────────────────────────────────────────────────────────────────────────────
\section{Conclusions and Future Work}

\subsection*{Summary of Findings}
This project demonstrates that a practical, reproducible pipeline can surface
longitudinal ethics signals from heterogeneous datasets. The Gemini keyword
corpus provides a usable proxy for ethics emphasis across years, while the
corporate policy timelines provide explicit policy points with category and
severity. Combined in a single UI, the system delivers a clear and professional
summary of policy evolution.

\subsection*{Significance and Impact}
The principal impact is a user-facing artefact that supports transparent
inspection of AI ethics commitments. The system is lightweight, deployable on
modest hardware, and designed for extension with additional datasets or richer
policy sources. It therefore serves as both a research tool and a practical
dashboard for stakeholders.

\subsection*{Limitations}
\begin{enumerate}
  \item The Gemini dataset is keyword-based and does not include full policy
  text; keyword counts should be treated as proxies rather than definitive
  statements.
  \item Corporate policy timelines cover a subset of companies and years;
  additional sources could improve longitudinal depth.
  \item Evaluation focuses on data coverage and interpretability rather than
  formal user studies; future work should include usability testing.
\end{enumerate}

\subsection*{Future Work}
\begin{enumerate}
  \item Integrate structured policy documents to align keyword signals with
  direct policy excerpts and citations.
  \item Extend the ratings system to include public feedback and inter-rater
  reliability metrics for policy summaries.
  \item Add advanced analytics (trend slopes, category deltas) and automated
  alerts for significant policy shifts.
\end{enumerate}

% ─────────────────────────────────────────────────────────────────────────────
\balance

\begin{thebibliography}{99}

\bibitem{jobin2019global}
A.~Jobin, M.~Ienca, and E.~Vayena,
``The global landscape of AI ethics guidelines,''
\textit{Nature Mach.\ Intell.}, vol.~1, no.~9, pp.~389--399, 2019.
\doi{10.1038/s42256-019-0088-2}

\bibitem{munn2023ethics}
L.~Munn,
``The ethics of AI ethics: An evaluation of guidelines,''
\textit{AI \& Ethics}, vol.~3, no.~2, pp.~1--12, 2023.
\doi{10.1007/s43681-022-00197-9}

\bibitem{springer2022development}
Springer Nature,
``Development of AI ethics guidelines model based on AI life cycle,''
\textit{Springer AI Ethics J.}, vol.~4, pp.~1--18, 2022.

\bibitem{stanford2023hai}
Stanford University HAI,
\textit{2023 AI Index Report},
Stanford, CA, USA, 2023.

\bibitem{nature2023responsible}
Nature Portfolio,
``Responsible AI measures dataset for ethics evaluation of AI systems,''
\textit{Scientific Data}, vol.~10, no.~1, pp.~1--9, 2023.

\bibitem{floridi2021ai}
L.~Floridi and J.~Cowls,
``A unified framework of five principles for AI in society,''
\textit{Philos.\ Trans.\ R.\ Soc.\ A}, vol.~379, no.~2207, 2021.
\doi{10.1098/rsta.2020.0367}

\bibitem{rahwan2021machine}
I.~Rahwan,
``Machine behaviour: A field of study to understand intelligent machines,''
\textit{Nature}, vol.~568, no.~7753, pp.~477--486, 2019.
\doi{10.1038/s41586-019-1138-y}

\bibitem{openai2023gpt4}
OpenAI,
``GPT-4 Technical Report,''
San Francisco, CA, USA, 2023, arXiv:2303.08774.

\bibitem{bommasani2022opportunities}
R.~Bommasani \textit{et~al.},
``On the opportunities and risks of foundation models,''
Stanford CRFM, Rep.\ 2022-001, 2022.

\bibitem{dataGov2024}
Data.gov,
``AI and machine learning datasets,''
U.S.\ GSA, Washington, D.C., USA, 2024.
[Online]. Available: \url{https://catalog.data.gov}

\bibitem{worldBank2024}
World Bank,
``World Bank Open Data,''
Washington, D.C., USA, 2024.
[Online]. Available: \url{https://data.worldbank.org}

\bibitem{ieeeEthics2024}
IEEE,
``Ethically Aligned Design (EAD v3),''
IEEE Standards Assoc., Piscataway, NJ, USA, 2024.

\bibitem{beijing2019ai}
Beijing Academy of Artificial Intelligence,
``Beijing AI Principles,''
2019. [Online]. Available: \url{https://www.baai.ac.cn}

\bibitem{eu2021ai}
European Commission,
``Ethics Guidelines for Trustworthy AI,''
Brussels, Belgium, 2021.

\bibitem{gebru2021datasheets}
T.~Gebru \textit{et~al.},
``Datasheets for datasets,''
\textit{Commun.\ ACM}, vol.~64, no.~12, pp.~86--92, 2021.

\end{thebibliography}

\end{document}
